\documentclass[11pt]{article}
\usepackage[a4paper,margin=1in]{geometry}
\usepackage{graphicx}
\usepackage{listings}
\usepackage{xcolor}
\usepackage{booktabs}
\usepackage{array}
\usepackage{hyperref}
\usepackage{parskip}
\hypersetup{colorlinks=true, linkcolor=black, urlcolor=blue}

\definecolor{codebg}{RGB}{245,245,245}
\lstdefinestyle{cppstyle}{
  language=C++,
  backgroundcolor=\color{codebg},
  basicstyle=\ttfamily\footnotesize,
  keywordstyle=\color{blue}\bfseries,
  commentstyle=\color{gray}\itshape,
  stringstyle=\color{teal},
  numbers=left,
  numberstyle=\tiny\color{gray},
  stepnumber=1,
  breaklines=true,
  frame=single,
  tabsize=2,
  showstringspaces=false
}
\lstdefinestyle{outstyle}{
  backgroundcolor=\color{codebg},
  basicstyle=\ttfamily\scriptsize,
  breaklines=true,
  frame=single,
  showstringspaces=false
}
\lstdefinestyle{bashstyle}{
  language=bash,
  backgroundcolor=\color{codebg},
  basicstyle=\ttfamily\footnotesize,
  breaklines=true,
  frame=single,
  showstringspaces=false
}

\title{\vspace{-2cm}}
\date{}

\begin{document}
\begin{titlepage}
\centering
\vspace*{1cm}
{\large \textbf{Sri Sivasubramaniya Nadar College of Engineering}}\\
{\small (An Autonomous Institution, Affiliated to Anna University, Chennai)}\\
{\small Kalavakkam -- 603 110}\\[0.3cm]
{\small Department of Computer Science \& Engineering}\\[1.5cm]

\rule{\linewidth}{0.5mm}\\[0.4cm]
{\LARGE \textbf{Computer Networks Laboratory}}\\[0.3cm]
{\Large 1513 -- Computer Networks Lab}\\[0.4cm]
\rule{\linewidth}{0.5mm}\\[1.5cm]

{\huge \textbf{Assignment 8}}\\[0.4cm]
{\Large Simulation and Performance Analysis of TCP and UDP Traffic Using NS-3}\\[2cm]

\begin{center}
\begin{tabular}{ll}
\textbf{Name:} & Sundareswaran Rajaguru \\[0.2cm]
\textbf{Register Number:} & 3122247001048 \\[0.2cm]
\textbf{Year / Semester:} & III Year / V Semester \\[0.2cm]
\textbf{Course:} & 1513 Computer Networks Lab \\[0.2cm]
\textbf{Date of Submission:} & 09 September 2026 \\[0.2cm]
\end{tabular}
\end{center}

\vfill
\end{titlepage}

\section*{1. Objective}
To develop an NS-3 program simulating packet transmission between four end
nodes through two intermediate routers, carrying one TCP flow and one UDP
flow across a shared bottleneck link, and to measure and compare packet
loss, average end-to-end delay, throughput, and TCP retransmission rate as
the bottleneck bandwidth, UDP source rate, queue size, and propagation
delay are independently varied.

\section*{2. Topology}
\begin{center}
\begin{verbatim}
n0 ---10Mbps/2ms--- r1 ===bottleneck=== r2 ---10Mbps/2ms--- n3
n1 ---10Mbps/2ms--- r1                  r2 ---10Mbps/2ms--- n2

TCP flow: n0 -> n3      UDP flow: n1 -> n2
\end{verbatim}
\end{center}
Access links (n0--r1, n1--r1, r2--n2, r2--n3) are fixed at 10~Mbps/2~ms. The
r1--r2 bottleneck link's bandwidth and propagation delay, its DropTail queue
size, and the UDP source rate are all command-line parameters, so a single
compiled program reproduces every run required below.

\section*{3. Algorithm}
\begin{enumerate}
  \item Create six nodes (\texttt{n0, n1, r1, r2, n2, n3}) and five
        point-to-point links: four fixed-rate access links plus the
        variable-rate bottleneck link between \texttt{r1} and \texttt{r2},
        whose queue is a \texttt{DropTailQueue} sized from the
        \texttt{--queueSize} argument.
  \item Install the Internet stack on all six nodes, assign one \texttt{/24}
        subnet per link, and populate global routing tables so traffic from
        \texttt{n0}/\texttt{n1} reaches \texttt{n3}/\texttt{n2} via
        \texttt{r1}$\to$\texttt{r2}.
  \item Fix the TCP segment size to 1024 bytes
        (\texttt{ns3::TcpSocket::SegmentSize}) and start a
        \texttt{BulkSendHelper} (continuous, congestion-controlled) TCP flow
        from \texttt{n0} to a \texttt{PacketSink} on \texttt{n3}.
  \item Start an \texttt{OnOffHelper} UDP flow (512-byte packets, rate from
        \texttt{--udpRate}, forced always-on) from \texttt{n1} to a
        \texttt{PacketSink} on \texttt{n2}.
  \item Enable pcap tracing on every point-to-point interface
        (\texttt{PointToPointHelper::EnablePcapAll}) to record packet-level
        transmission/reception, and install \texttt{FlowMonitor} on all
        nodes to classify and count packets per flow.
  \item Run the simulation for 20 seconds, then read back per-flow
        transmitted/received/lost packet counts and cumulative delay from
        \texttt{FlowMonitor} and derive packet loss~\%, throughput, average
        delay, and (for the TCP flow) an approximate retransmission rate.
\end{enumerate}

\section*{4. Code}
\subsection*{4.1 ass8.cc}
\lstinputlisting[style=cppstyle]{ns-3-work/scratch/ass8.cc}

\section*{5. Baseline Output}
Run with the spec's exact defaults: bottleneck 2~Mbps/10~ms, DropTail queue
of 50 packets, UDP rate 1~Mbps, 20-second simulation.
\begin{lstlisting}[style=outstyle]
[TCP] 10.0.1.1:49153 -> 10.0.5.2:9001 | Tx=2254 Rx=2187 Lost=44 (1.95%) | Throughput=0.990Mbps | AvgDelay=173.6ms | RetransRate=2.97%
[UDP] 10.0.2.1:49153 -> 10.0.4.2:9002 | Tx=4638 Rx=4403 Lost=196 (4.23%) | Throughput=1.001Mbps | AvgDelay=170.8ms
[TCP] 10.0.5.2:9001 -> 10.0.1.1:49153 | Tx=1296 Rx=1295 Lost=0 (0.00%) | Throughput=0.031Mbps | AvgDelay=14.3ms  | RetransRate=0.08%
\end{lstlisting}
The third line is the TCP connection's reverse-direction ACK traffic
(\texttt{n3}$\to$\texttt{n0}), also captured by \texttt{FlowMonitor} --
expected for any full-duplex TCP flow, not a separate experiment.

A packet-level trace confirms both protocols on the wire exactly as
configured (TCP three-way handshake followed by 1024-byte data segments):
\begin{lstlisting}[style=outstyle]
00:00:01.002046 IP 10.0.1.1.49153 > 10.0.5.2.9001: Flags [S], seq 0, win 65535, ...
00:00:01.106603 IP 10.0.5.2.9001 > 10.0.1.1.49153: Flags [S.], seq 0, ack 1, win 65535, ...
00:00:01.110692 IP 10.0.1.1.49153 > 10.0.5.2.9001: Flags [.], ack 1, win 32768, ...
00:00:01.111555 IP 10.0.1.1.49153 > 10.0.5.2.9001: Flags [.], seq 1:1025, ack 1, ..., length 1024
00:00:01.112417 IP 10.0.1.1.49153 > 10.0.5.2.9001: Flags [.], seq 1025:2049, ack 1, ..., length 1024
\end{lstlisting}

\section*{6. Full Output (All Experiments)}
\lstinputlisting[style=outstyle]{output.txt}

\section*{7. Experiment 1 -- Bottleneck Bandwidth}
\begin{center}
\renewcommand{\arraystretch}{1.2}
\begin{tabular}{lrrrr}
\toprule
\textbf{Bottleneck} & \textbf{TCP Loss \%} & \textbf{TCP Throughput (Mbps)} & \textbf{Avg Delay (ms)} & \textbf{TCP Retrans \%} \\
\midrule
2 Mbps  & 1.95 & 0.990 & 173.6 & 2.97 \\
5 Mbps  & 0.21 & 3.927 & 83.9  & 0.57 \\
10 Mbps & 0.05 & 8.890 & 54.2  & 0.36 \\
\bottomrule
\end{tabular}
\end{center}
As bottleneck bandwidth increases, TCP's congestion window can grow much
larger before hitting a loss event, so throughput rises roughly in
proportion to capacity, while both loss and (crucially) queuing delay fall
sharply -- at 2~Mbps the queue is almost always close to full (173~ms of
delay, dominated by queuing, not the 10~ms link propagation delay), whereas
at 10~Mbps the same offered load barely troubles the queue at all.

\section*{8. Experiment 2 (Optional) -- UDP Traffic Rate}
\begin{center}
\renewcommand{\arraystretch}{1.2}
\begin{tabular}{lrrrr}
\toprule
\textbf{UDP Rate} & \textbf{UDP Loss \%} & \textbf{UDP Delay (ms)} & \textbf{TCP Throughput (Mbps)} & \textbf{TCP Retrans \%} \\
\midrule
500 Kbps & 0.00  & 174.4  & 1.469 & 1.55 \\
1 Mbps   & 4.23  & 170.8  & 0.990 & 2.97 \\
2 Mbps   & 46.77 & 658.2  & 0.935 & 2.22 \\
5 Mbps   & 38.98 & 4926.2 & 0.930 & 2.10 \\
\bottomrule
\end{tabular}
\end{center}
This is the clearest illustration of TCP vs.\ UDP fairness: since the
bottleneck is only 2~Mbps, a 500~Kbps--1~Mbps UDP flow fits comfortably and
TCP still gets most of the remaining capacity. Once UDP is offered above the
bottleneck's capacity (2~Mbps, 5~Mbps), UDP itself suffers severe loss
(up to 46.8\%) and its own queuing delay explodes into seconds, because UDP
never slows down in response -- yet TCP's throughput barely drops
(0.99~Mbps~$\to$~0.93~Mbps), because TCP's congestion control keeps
backing off just enough to preserve its fair share rather than collapsing.
UDP is happy to hurt itself (and, to a lesser extent, TCP) rather than yield
bandwidth voluntarily.

\section*{9. Experiment 3 (Optional) -- Queue Size}
\begin{center}
\renewcommand{\arraystretch}{1.2}
\begin{tabular}{lrrr}
\toprule
\textbf{Queue Size} & \textbf{TCP Loss \%} & \textbf{Avg Delay (ms)} & \textbf{TCP Retrans \%} \\
\midrule
20 packets  & 2.55 & 92.9  & 3.00 \\
50 packets  & 1.95 & 173.6 & 2.97 \\
100 packets & 1.32 & 298.8 & 2.59 \\
\bottomrule
\end{tabular}
\end{center}
This is a textbook \emph{bufferbloat} trade-off: a larger queue absorbs more
of TCP's burst without dropping packets (loss falls from 2.55\% to 1.32\%),
but every packet that would otherwise have been dropped instead sits longer
in the queue, so average end-to-end delay more than triples (93~ms to
299~ms) as the queue grows from 20 to 100 packets. A bigger buffer does not
make the link faster; it just trades packet loss for latency.

\section*{10. Experiment 4 (Optional) -- Propagation Delay}
\begin{center}
\renewcommand{\arraystretch}{1.2}
\begin{tabular}{lrrr}
\toprule
\textbf{Bottleneck Delay} & \textbf{TCP Loss \%} & \textbf{Avg Delay (ms)} & \textbf{TCP Retrans \%} \\
\midrule
5 ms  & 0.95 & 163.1 & 1.86 \\
10 ms & 1.95 & 173.6 & 2.97 \\
20 ms & 1.97 & 181.4 & 2.97 \\
50 ms & 3.83 & 204.4 & 4.90 \\
\bottomrule
\end{tabular}
\end{center}
Average end-to-end delay rises only modestly with propagation delay because
queuing delay (governed by bandwidth and queue size, not propagation delay)
still dominates the total at every setting. However, a longer round-trip
time slows down how quickly TCP's loss-based congestion control can react
to congestion signals (ACKs take longer to return), so both packet loss and
retransmission rate climb noticeably at 50~ms as TCP overshoots the
available buffer more before backing off.

\section*{11. Overall Comparison of TCP and UDP}
Across every experiment, TCP consistently adapts its sending rate to fit
whatever capacity remains after UDP, at the cost of retransmitting the small
fraction of segments that are lost before it reacts; UDP never adapts, so
when it is offered more than the bottleneck can carry it is UDP's own
packets (not TCP's) that are dropped in bulk and delayed severely. This
matches the standard textbook result: TCP protects the network (and other
flows) from itself, while an unconstrained UDP source can degrade its own
delivered quality without ever being told to slow down.

\section*{12. Learning Outcomes}
\begin{itemize}
  \item Built a dumbbell (bottleneck) topology in NS-3 with separate access
        and bottleneck link characteristics, and understood why this shape
        is the standard setup for studying congestion.
  \item Used \texttt{FlowMonitor} to obtain per-flow packet loss, delay, and
        throughput statistics without hand-writing trace callbacks, and
        derived an approximate TCP retransmission rate from its Tx/Rx packet
        counts.
  \item Observed experimentally that TCP throughput scales with bottleneck
        bandwidth while queuing delay falls, and that a larger DropTail
        queue trades packet loss for increased latency (bufferbloat).
  \item Demonstrated concretely why sending UDP above the bottleneck's
        capacity harms the UDP flow itself far more than the competing TCP
        flow, since UDP has no congestion-avoidance feedback loop.
  \item Practised enabling NS-3's built-in pcap tracing
        (\texttt{EnablePcapAll}) to verify, at the packet level in
        Wireshark/tcpdump, that the TCP handshake, segment size, and UDP
        packet size all matched the configured parameters.
  \item Practised parameterising a single NS-3 program via
        \texttt{CommandLine} so that all four required experiments (varying
        bandwidth, UDP rate, queue size, propagation delay) are produced by
        one binary and a short shell sweep, without recompiling per data
        point.
\end{itemize}

\end{document}
